\chapter{2023年上海卷（秋）}

\section{填空题}

\begin{problem}
设 \(x\in\R\)，不等式 \(\left|x-2\right|<1\) 的解集为\fillinblank{}.
\end{problem}

\begin{answer}
\(\left(1,3\right)\)
\end{answer}

\begin{solution}
由 \(\left|x-2\right|<1\) 得 \(-1<x-2<1\)，所以 \(1<x<3\). 故解集为 \(\left(1,3\right)\).
\end{solution}

\begin{problem}
若向量 \(\bs a=\left(-2,3\right)\)，\(\bs b=\left(1,2\right)\)，则 \(\bs a\cdot\bs b=\fillinblank{}\).
\end{problem}

\begin{answer}
\(4\)
\end{answer}

\begin{solution}
\(\bs a\cdot\bs b=-2\times 1+3\times 2=4\).
\end{solution}

\begin{problem}
若数列 \(\{a_n\}\) 是首项为 \(3\)，公比为 \(2\) 的等比数列，其前 \(n\) 项和为 \(S_n\)，则 \(S_6=\fillinblank{}\).
\end{problem}

\begin{answer}
\(189\)
\end{answer}

\begin{solution}
\(S_6=3\frac{1-2^6}{1-2}=3\times 63=189\)
\end{solution}

\begin{problem}
若 \(\tan A=3\)，则 \(\tan 2A=\fillinblank{}\).
\end{problem}

\begin{answer}
\(-\frac34\)
\end{answer}

\begin{solution}
\[
\tan 2A=\frac{2\tan A}{1-\tan^2 A}
=\frac{2\times 3}{1-3^2}
=-\frac34
\]
\end{solution}

\begin{problem}
函数 \(y=\begin{cases}
2^x,&x>0,\\
1,&x\le 0,
\end{cases}\) 的值域为\fillinblank{}.
\end{problem}

\begin{answer}
\(\left[1,+\infty\right)\)
\end{answer}

\begin{solution}
当 \(x\le 0\) 时，\(y=1\)；当 \(x>0\) 时，\(y=2^x>1\). 故值域为 \(\left[1,+\infty\right)\).
\end{solution}

\begin{problem}
若复数 \(z=1+\i\)（\(\i\) 为虚数单位），则 \(\left|1-\i z\right|=\fillinblank{}\).
\end{problem}

\begin{answer}
\(\sqrt5\)
\end{answer}

\begin{solution}
\(1-\i z=1-\i(1+\i)=2-\i\)，所以 \(\left|1-\i z\right|=\sqrt{2^2+(-1)^2}=\sqrt5\).
\end{solution}

\begin{problem}
设 \(m\in\R\). 若圆 \(x^2+y^2-4y-m=0\) 的面积为 \(\pi\)，则 \(m=\fillinblank{}\).
\end{problem}

\begin{answer}
\(-3\)
\end{answer}

\begin{solution}
\(x^2+y^2-4y-m=0\iff x^2+\left(y-2\right)^2=m+4\)，圆的面积为 \(\pi\)，故半径为 \(1\)，于是 \(m+4=1\)，解得 \(m=-3\).
\end{solution}

\begin{problem}
在 \(\triangle ABC\) 中，角 \(A\)，\(B\)，\(C\) 所对的边长分别为 \(a\)，\(b\)，\(c\). 若 \(a=4\)，\(b=5\)，\(c=6\)，则 \(\sin A=\fillinblank{}\).
\end{problem}

\begin{answer}
\(\frac{\sqrt7}{4}\)
\end{answer}

\begin{solution}
由余弦定理，
\[
\cos A=\frac{b^2+c^2-a^2}{2bc}
=\frac{25+36-16}{2\times 5\times 6}
=\frac34
\]
故
\[
\sin A=\sqrt{1-\cos^2 A}=\sqrt{1-\frac{9}{16}}=\frac{\sqrt7}{4}
\]
\end{solution}

\begin{problem}
国内生产总值（GDP）是衡量地区经济状况的最佳指标. 根据统计数据显示，某市在 2020 年间经济高质量增长，GDP 稳步增长，第一季度和第四季度的 GDP 分别为 \(232\) 亿元和 \(241\) 亿元，且四个季度 GDP 逐季度增长. 若这四个季度 GDP 的中位数与平均数相等，则该市 2020 年的 GDP 总额为\fillinblank{} 亿元.
\end{problem}

\begin{answer}
\(946\)
\end{answer}

\begin{solution}
设第二，三季度 GDP 分别为 \(b\)，\(c\). 因逐季度增长，故 \(232\le b\le c\le 241\). 中位数与平均数相等，即
\(\frac{b+c}{2}=\frac{232+b+c+241}{4}\)，化简得 \(b+c=473\)，故全年 GDP 总额为
\(232+b+c+241=232+473+241=946\)
\end{solution}

\begin{problem}
已知对任意给定的实数 \(x\)，都有
\(
\left(1+2023x\right)^{100}+\left(2023-x\right)^{100}=a_0+a_1x+a_2x^2+\cdots+a_{100}x^{100},
\)
其中 \(a_0\)，\(a_1\)，\(\cdots\)，\(a_{100}\in\R\). 若 \(a_k<0\)，其中 \(k\in\{0,1,2,\cdots,100\}\)，则正整数 \(k\) 的最大值为\fillinblank{}.
\end{problem}

\begin{answer}
\(49\)
\end{answer}

\begin{solution}
展开得 \(a_k=\mathrm C_{100}^{k}2023^k+\mathrm C_{100}^{k}(-1)^k2023^{100-k}=\mathrm C_{100}^{k}\bigl(2023^k+(-1)^k2023^{100-k}\bigr)\)，当 \(k\) 为偶数时，\(a_k>0\). 当 \(k\) 为奇数时，
\(a_k<0\iff 2023^k<2023^{100-k}\iff k<50\)，
故正整数 \(k\) 的最大值为 \(49\).
\end{solution}

\begin{problem}
某公园欲建设一段斜坡，斜坡面是一平面，坡面与水平地面所成角为 \(\theta\)，坡顶距离水平地面的垂直高度为 \(4\) 米，游客从坡底沿着斜坡面向上直行，每走 \(1\) 米消耗的体力为 \(1.025-\cos\theta\)，当 \(\theta=\fillinblank{}\) 时，游客走斜坡所消耗的总体力最小.
\end{problem}

\begin{answer}
\(\arccos\frac{40}{41}\)
\end{answer}

\begin{solution}
斜坡长为 \(\frac{4}{\sin\theta}\)，故总消耗体力 \(E(\theta)=\frac{4\left(1.025-\cos\theta\right)}{\sin\theta}\)，求导得
\(E'(\theta)=\frac{4\left(1-1.025\cos\theta\right)}{\sin^2\theta}\)，令 \(E'(\theta)=0\)，得 \(\cos\theta=\frac{40}{41}\). 又 \(E'(\theta)\) 在该点左右变号，故当 \(\theta=\arccos\frac{40}{41}\) 时，总体力最小.
\end{solution}

\begin{problem}
空间中有三个定点 A，B，C，且 \(AB=BC=CA=1\). 若在空间中任取 \(2\) 个不同的点，使它们与 A，B，C 恰好成为一个正四棱锥的五个顶点，则不同的取法共有\fillinblank{} 种.
\end{problem}

\begin{answer}
\(9\)
\end{answer}

\begin{solution}
分两类讨论.
当 \(\triangle ABC\) 为正四棱锥的侧面时，平面 \(ABC\) 的两侧分别可以做 \(ABDE\) 作为底面，如图有 \(2\) 种情况；分别以 \(BCED\)，\(ACED\) 为底面也各有 \(2\) 种情况，所以共有 \(6\) 种情况.

当 \(\triangle ABC\) 为正四棱锥的截面时，如图，\(D\)，\(E\) 位于 \(AB\) 两侧，\(ADBE\) 为底面，只有 \(1\) 种情况；分别以 \(BDCE\)，\(ADCE\) 为底面也各有 \(1\) 种情况，所以共有 \(3\) 种情况.

综上，共有 \(6+3=9\) 种情况.
\end{solution}
\section{单选题}

\begin{problem}
已知集合 \(P=\{1,2\}\)，\(Q=\{2,3\}\). 若 \(M=\{x\mid x\in P,\ x\notin Q\}\)，则 \(M=\quad\).

\choices
    {\(\{1\}\)}
    {\(\{2\}\)}
    {\(\{3\}\)}
    {\(\{1,2,3\}\)}
\end{problem}

\begin{answer}
A
\end{answer}

\begin{solution}
由定义知 \(M=\{1\}\)，故选 A.
\end{solution}

\begin{problem}
为了解某校高中男生身高与体重的关系，随机选取 \(50\) 名男生的身高与体重的数据，绘制散点图，如图所示，下列说法正确的是（\quad）.

\bitmapfigure[max width=0.68\linewidth]{img/2023/shanghai/q14_fig1.png}

\choices
    {身高越高，体重越重}
    {身高越高，体重越轻}
    {身高与体重之间呈正相关}
    {身高与体重之间呈负相关}
\end{problem}

\begin{answer}
C
\end{answer}

\begin{solution}
由散点图分布可知，身高与体重成正相关，故选 C.
\end{solution}

\begin{problem}
已知 \(a>0\)，函数 \(y=\sin x\) 在区间 \(\left[a,2a\right]\) 上的最小值为 \(s\)，在区间 \(\left[2a,3a\right]\) 上的最小值为 \(t\)，当 \(a\) 变化时，下列一定不成立的是（\quad）.

\choices
    {\(s>0,t>0\)}
    {\(s>0\)，\(t<0\)}
    {\(s<0\)，\(t<0\)}
    {\(s<0\)，\(t>0\)}
\end{problem}

\begin{answer}
D
\end{answer}

\begin{solution}
若 \(t>0\)，则区间 \(\left[2a,3a\right]\) 必须完全落在某个正值区间 \(\left(2k\pi,(2k+1)\pi\right)\) 内.
当 \(k=0\) 时，\(3a<\pi\)，从而 \(a<\frac{\pi}{3}\)，于是 \(\left[a,2a\right]\subset(0,\pi)\)，故 \(s>0\)；当 \(k\ge 1\) 时，\(2a>2k\pi\) 与 \(3a<(2k+1)\pi\) 不能同时成立.
因此 \(s<0\)，\(t>0\) 不可能，故选 D.
\end{solution}

\begin{problem}
对于曲线 \(C\)，若存在点 \(M\)，使得对任意的点 \(P\in C\)，都存在 \(Q\in C\)，且 \(\left|PM\right|\cdot\left|QM\right|=1\)，则称曲线 \(C\) 为“自相关曲线”. 有以下结论：
\begin{circlelist}
\item 任何椭圆都是“自相关曲线”；
\item 存在双曲线是“自相关曲线”. 关于上述两个结论，说法正确的是（\quad）.
\end{circlelist}

\choices
    {\circled{1} 成立，\circled{2} 成立}
    {\circled{1} 成立，\circled{2} 不成立}
    {\circled{1} 不成立，\circled{2} 成立}
    {\circled{1} 不成立，\circled{2} 不成立}
\end{problem}

\begin{answer}
B
\end{answer}

\begin{solution}
\circled{1} 成立. 设 \(d_-(M)=\min_{P\in C}|PM|\)，\(d_+(M)=\max_{P\in C}|PM|\). 椭圆 \(C\) 是紧且连通的，故距离集合为区间 \([d_-(M),d_+(M)]\). 当 \(M\) 取在椭圆上时 \(d_-(M)d_+(M)=0\)；当 \(M\) 远离椭圆时 \(d_-(M)d_+(M)\to+\infty\). 由连续性可取某个 \(M\) 使 \(d_-(M)d_+(M)=1\)，此时区间为 \([d_-(M),1/d_-(M)]\)，对倒数运算封闭，故任何椭圆都是"自相关曲线".

\circled{2} 不成立. 双曲线是闭集且无界。若固定点 \(M\notin C\)，则 \(d_-(M)>0\)，而距离集合含有任意大的数；取距离足够大的点 \(P\)，则 \(1/|PM|<d_-(M)\)，不可能找到 \(Q\in C\) 满足 \(|PM||QM|=1\). 若 \(M\in C\)，取 \(P=M\) 时左边恒为 \(0\)，同样不可能。因此不存在双曲线是"自相关曲线". 故选 B.
\end{solution}

\section{解答题}

\begin{problem}
2023 年 6 月 7 日，21 世纪汽车博览会在上海举行. 已知某汽车模型公司共有 \(25\) 个汽车模型，其外观和内饰的颜色分布如下表所示：

\begin{center}
\begin{tabular}{c|cc}
\hline
& 红色外观 & 蓝色外观\\
\hline
米色内饰 & \(12\) & \(8\)\\
棕色内饰 & \(2\) & \(3\)\\
\hline
\end{tabular}
\end{center}

\begin{enumerate}
\item 小明从这些汽车模型中随机取出一个，记事件 A 为“小明取到红色外观的汽车模型”，事件 B 为“小明取到棕色内饰的汽车模型”，求 \(P(B)\) 和 \(P(B\mid A)\)，并据此判断事件 A 和事件 B 是否独立；
\item 该公司举行抽奖活动，规定在一次抽奖中，每人可以一次性从这些汽车模型中随机取出两个汽车模型，并给出以下假设：
\begin{enumerate}
\item 取出的两个汽车模型会出现三种结果，即外观和内饰均为同色，外观和内饰均不为同色，以及仅外观或仅内饰为同色；
\item 按出现结果的可能性大小，概率越小奖项越高；
\item 抽奖活动的奖金金额为：一等奖 \(600\) 元，二等奖 \(300\) 元，三等奖 \(150\) 元.
\end{enumerate}
请分析各奖项对应的结果，设 \(X\) 为奖金额，写出 \(X\) 的分布并求其数学期望.
\end{enumerate}
\end{problem}

\begin{solution}
\begin{enumerate}
\item 红色外观共有 \(12+2=14\) 个，棕色内饰共有 \(2+3=5\) 个，所以 \(P(B)=\frac{5}{25}=\frac15\)，\(P(B\mid A)=\frac{2}{14}=\frac17\).
因为 \(P(B\mid A)\ne P(B)\)，所以事件 \(A\) 和事件 \(B\) 不独立.

\item 从 \(25\) 个模型中一次取出 \(2\) 个，不计顺序，总取法数为 \(\mathrm C_{25}^{2}=300\).
“外观和内饰均为同色”即两模型完全相同，取法数为 \(\mathrm C_{12}^{2}+\mathrm C_{8}^{2}+\mathrm C_{2}^{2}+\mathrm C_{3}^{2}=98\)，故其概率为 \(\frac{98}{300}=\frac{49}{150}\).

“仅外观或仅内饰为同色”的取法数为 \(12\times 2+8\times 3+12\times 8+2\times 3=150\)，故其概率为 \(\frac{150}{300}=\frac12\).

“外观和内饰均不为同色”的取法数为 \(12\times 3+8\times 2=52\)，故其概率为 \(\frac{52}{300}=\frac{13}{75}\).

由“概率越小奖项越高”可知：\(P(X=600)=\frac{13}{75}\)，\(P(X=300)=\frac{49}{150}\)，\(P(X=150)=\frac12\).
因此 \(E(X)=600\times \frac{13}{75}+300\times \frac{49}{150}+150\times \frac12=277\).
故 \(X\) 的分布列为
{\centering
\begin{tabular}{c|ccc}
\(X\) & \(600\) & \(300\) & \(150\) \\
\hline
\(P\) & \(\frac{13}{75}\) & \(\frac{49}{150}\) & \(\frac12\)
\end{tabular}
\par}
且 \(E(X)=277\).
\end{enumerate}
\end{solution}

\begin{problem}
在四棱柱 \(ABCD-A_1B_1C_1D_1\) 中，\(AB\parallel DC\)，\(AB\perp AD\)，\(AB=2\)，\(AD=3\)，\(DC=4\).

\bitmapfigure[max width=0.36\linewidth]{img/2023/shanghai/q17_fig1.png}
\end{problem}

\begin{solution}
在四棱柱中，\(AA_1\parallel DD_1\)，且 \(AB\parallel DC\). 因此 \(\overrightarrow{A_1B}=\overrightarrow{AB}-\overrightarrow{AA_1}\) 的方向可由平面 \(DCC_1D_1\) 内的两个方向 \(\overrightarrow{DC}\)，\(\overrightarrow{DD_1}\) 线性表示，故 \(A_1B\parallel\) 平面 \(DCC_1D_1\).

取坐标系：令 \(A(0,0,0)\)，\(B(2,0,0)\)，\(D(0,3,0)\)，\(C(4,3,0)\). 直四棱柱的高为 \(h\)，则 \(A_1(0,0,h)\)，\(B_1(2,0,h)\)，\(C_1(4,3,h)\)，\(D_1(0,3,h)\). 由棱柱体积为 \(36\) 得底面四边形 \(ABCD\) 的面积 \(S_{ABCD}=\frac{2+4}{2}\times 3=9\)，故棱柱高 \(h=\frac{36}{9}=4\).
设 \(H\) 为点 \(A\) 到直线 \(BD\) 的垂足，则在直角三角形 \(ABD\) 中，\(BD=\sqrt{2^2+3^2}=\sqrt{13}\)，\(AH=\frac{AB\cdot AD}{BD}=\frac{6}{\sqrt{13}}\). 二面角 \(A_1-BD-A\) 的平面角就是 \(\angle A_1HA\)，因此 \(\tan\angle A_1-BD-A=\frac{AA_1}{AH}=\frac{4}{6/\sqrt{13}}=\frac{2\sqrt{13}}{3}\). 故 \(\angle A_1-BD-A=\arctan\frac{2\sqrt{13}}{3}\)，也可写成 \(\arccos\frac{3}{\sqrt{61}}\).
\end{solution}

\begin{problem}
设函数 \(f(x)=\frac{x^2+(3a+1)x+c}{x+a}\)，\(a\)，\(c\in\R\)
\begin{enumerate}
\item 当 \(a=0\) 时，是否存在实数 \(c\)，使得 \(y=f(x)\) 为奇函数，请说明理由；
\item 若函数 \(y=f(x)\) 的图像过点 \((1,3)\)，且其与 \(x\) 轴的负半轴有两个不同的交点，求 \(c\) 的值及 \(a\) 的取值范围.
\end{enumerate}
\end{problem}

\begin{solution}
\begin{enumerate}
\item 当 \(a=0\) 时，\(f(x)=\frac{x^2+x+c}{x}=x+1+\frac{c}{x}\)，于是对任意 \(x\ne 0\)，有 \(f(-x)=-x+1-\frac{c}{x}\)，\(f(x)+f(-x)=2\neq 0\)，所以不存在这样的实数 \(c\).

\item 由图像过点 \((1,3)\)，得 \(f(1)=3\iff \frac{1+(3a+1)+c}{1+a}=3\)，化简得 \(c=1\).
此时 \(f(x)=\frac{x^2+(3a+1)x+1}{x+a}\)，函数图像与 \(x\) 轴的交点由分子 \(x^2+(3a+1)x+1=0\) 给出. 要使其在负半轴上有两个不同交点，需要这两个根都为负且不同.
因为两根之积为 \(1>0\)，故只要两根和为负且判别式大于 \(0\) 即可：\(-(3a+1)<0\)，\((3a+1)^2-4>0\).
前者给出 \(a>-\frac13\)，后者在此条件下化为 \(3a+1>2\)，即 \(a>\frac13\).
此外，当 \(a=\frac12\) 时，\(x^2+\left(3a+1\right)x+1=x^2+\frac52x+1=\left(x+\frac12\right)\left(x+2\right)\)，与分母 \(x+a=x+\frac12\) 约去后，只剩一个负半轴交点，因此还需排除 \(a=\frac12\).
故 \(c=1\)，\(a\in\left(\frac13,+\infty\right)\setminus\left\{\frac12\right\}\).
\end{enumerate}
\end{solution}

\begin{problem}
已知抛物线 \(\Gamma:y^2=4x\)，第一象限内点 \(A\) 在 \(\Gamma\) 上，\(A\) 的纵坐标是 \(a\).
\begin{enumerate}
\item 若 \(A\) 到准线的距离为 \(3\)，求 \(a\)；
\item 若 \(a=4\)，\(B\) 在 \(x\) 轴上，且线段 \(AB\) 的中点在 \(\Gamma\) 上，求点 \(B\) 的坐标及坐标原点 \(O\) 到直线 \(AB\) 的距离；
\item 直线 \(l:x=-3\)，令 \(P\) 是第一象限 \(\Gamma\) 上异于 A 的一点，直线 \(AP\) 交 \(l\) 于 \(Q\)，\(H\) 是 \(P\) 在 \(l\) 上的投影，若点 A 满足“对于任意 \(P\) 都有 \(\left|HQ\right|>4\)”，求 \(a\) 的取值范围.
\end{enumerate}
\end{problem}

\begin{solution}
\begin{enumerate}
\item 设点 \(A\left(\frac{a^2}{4},a\right)\). 抛物线 \(y^2=4x\) 的准线为 \(x=-1\)，点 \(A\) 到准线的距离为 \(\frac{a^2}{4}+1\). 由题意 \(\frac{a^2}{4}+1=3\)，得 \(a=2\sqrt2\).

\item 当 \(a=4\) 时，\(A(4,4)\). 设 \(B(x,0)\)，则 \(AB\) 的中点为 \(\left(\frac{x+4}{2},2\right)\). 该点在 \(\Gamma\) 上，所以 \(2^2=4\cdot \frac{x+4}{2}\)，解得 \(x=-2\)，故 \(B(-2,0)\). 直线 \(AB\) 的方程为 \(2x-3y+4=0\)，于是原点 \(O\) 到直线 \(AB\) 的距离为 \(\frac{4}{\sqrt{2^2+(-3)^2}}=\frac{4}{\sqrt{13}}\).

\item 设 \(P\left(\frac{u^2}{4},u\right)\)，其中 \(u>0\) 且 \(u\ne a\). 又 \(A\left(\frac{a^2}{4},a\right)\)，则直线 \(AP\) 的斜率为 \(\frac{u-a}{\frac{u^2-a^2}{4}}=\frac{4}{u+a}\). 由此得直线 \(AP\) 与 \(x=-3\) 交于 \(Q\left(-3,\frac{au-12}{u+a}\right)\). 而 \(H(-3,u)\)，故 \(\left|HQ\right|=u-\frac{au-12}{u+a}=\frac{u^2+12}{u+a}\).
题设要求 \(\frac{u^2+12}{u+a}>4\) 对任意 \(u>0\)，\(u\ne a\) 成立，即 \(u^2-4u+12-4a>0\).
左边视为关于 \(u\) 的二次函数，其最小值在 \(u=2\) 处取得，最小值为 \(8-4a\).
因此若 \(a<2\)，则不等式恒成立；若 \(a=2\)，则左边化为 \((u-2)^2\)，其等号只在 \(u=2\) 处取得，而这时 \(P=A\) 被排除，所以仍成立；若 \(a>2\)，取 \(u=2\) 即不成立.
又 \(A\) 在第一象限，故 \(a>0\).
所以 \(a\in\left(0,2\right]\).
\end{enumerate}
\end{solution}

\begin{problem}
已知 \(f(x)=\ln x\). 曲线 \(y=f(x)\) 在点 \(\left(a_1,f(a_1)\right)\) 处的切线交 \(y\) 轴于点 \((0,a_2)\)，且 \(a_2>0\)；曲线 \(y=f(x)\) 在点 \(\left(a_2,f(a_2)\right)\) 处的切线交 \(y\) 轴于点 \((0,a_3)\)；以此类推，直至 \(a_m\le 0\) 时停止，得到数列 \(\{a_n\}\).
\begin{enumerate}
\item 证明：当正整数 \(n\ge 2\) 时，\(a_n=\ln a_{n-1}-1\)；
\item 当正整数 \(n\ge 2\) 时，试比较 \(a_n\) 与 \(a_{n-1}-2\) 的大小；
\item 若正整数 \(k\ge 3\)，是否存在 \(k\) 使得 \(a_1\)，\(a_2\)，\(\cdots\)，\(a_k\) 依次成等差数列？若存在，求出 \(k\) 的所有取值；若不存在，试说明理由.
\end{enumerate}
\end{problem}

\begin{solution}
\begin{enumerate}
\item 曲线 \(y=\ln x\) 在点 \(\left(a_{n-1},\ln a_{n-1}\right)\) 处的切线方程为 \(y-\ln a_{n-1}=\frac1{a_{n-1}}\left(x-a_{n-1}\right)\). 令 \(x=0\)，得该切线与 \(y\) 轴的交点纵坐标为 \(\ln a_{n-1}-1\)，所以 \(a_n=\ln a_{n-1}-1\).

\item 由 \(\ln x\le x-1\)（\(x>0\)）得 \(a_n=\ln a_{n-1}-1\le a_{n-1}-2\)，当且仅当 \(a_{n-1}=1\) 时取等号.

\item 记 \(h(x)=e^{x+1}+\ln x-1-2x\)（\(x>0\)）. 若 \(a_1\)，\(a_2\)，\(a_3\) 成等差数列，则 \(a_1+a_3=2a_2\). 由 （1） 知 \(a_1=e^{a_2+1}\)，且 \(a_3=\ln a_2-1\)，故 \(h(a_2)=0\). 又 \(h'(x)=e^{x+1}+\frac1x-2>0\)（\(x>0\)），所以 \(h(x)\) 在 \((0,+\infty)\) 上严格递增. 并且 \(\lim\limits_{x\to 0^+}h(x)=-\infty\)，\(h(1)=e^2-3>0\)，故方程 \(h(x)=0\) 在 \((0,1)\) 内有唯一解，从而存在唯一的 \(a_2\in(0,1)\) 使得 \(a_1\)，\(a_2\)，\(a_3\) 成等差数列. 此时 \(a_3=\ln a_2-1<0\)，所以序列在 \(a_3\) 处停止，满足题意的 \(k\) 只能取 \(3\).
\end{enumerate}
\end{solution}
